\newpage\setcounter{dang}{0}
%----------------NỘI DUNG BIÊN SOẠN
% \setcounter{section}{13}
\section{Các số đặc trưng đo mức độ phân tán}
%-------------------------------------------
\subsection{KIẾN THỨC TRỌNG TÂM}
\subsubsection{Khoảng biến thiên và khoảng tứ phân vị}
\begin{boxdn}
	Khoảng biến thiên, kí hiệu là $R$, là hiệu số giữa giá trị lớn nhất và giá trị nhỏ nhất trong mẫu số liệu.
\end{boxdn}
\textbf{Ý nghĩa.} Khoảng biến thiên dùng để đo độ phân tán của mẫu số liệu. Khoảng biến thiên càng lớn thì mẫu số liệu càng phân tán.
\begin{nx}
Sử dụng khoảng biến thiên có ưu điểm là đơn giản, dễ tính toán song khoảng biến thiên chỉ sử dụng thông tin của giá trị lớn nhất và giá trị nhỏ nhất mà bỏ qua thông tin từ tất cả các giá trị khác. Do đó, khoảng biến thiên rất dễ bị ảnh hưởng bởi các giá trị bất thường.
\end{nx}
\begin{boxdn}
	Khoảng tứ phân vị, kí hiệu là $\Delta_Q$, là hiệu số giữa tứ phân vị thứ ba và tứ phân vị thứ nhất, tức là $$\Delta_Q = Q_3 - Q_1.$$
\end{boxdn}
\subsubsection{Phương sai và độ lệch chuẩn}
\begin{boxdn}
	Cho mẫu số liệu $x_1, x_2, \dots, x_n$, nếu gọi số trung bình là $\overline{x}$ thì với mỗi giá trị $x_i$, độ lệch của nó so với giá trị trung bình là $x_i - \overline{x}$.
		\begin{itemize}
			\item Phương sai là giá trị
			$s^2 = \dfrac{(x_1 - \overline{x})^2 + (x_2 - \overline{x})^2 + \dots + (x_n - \overline{x})^2}{n}$.
			\item Căn bậc hai của phương sai, $s = \sqrt{s^2}$, được gọi là độ lệch chuẩn.
		\end{itemize}
\end{boxdn}
% \begin{note}
% 	Người ta còn sử dụng đại lượng để đo độ phân tán của mẫu số liệu:
% 	\[
% 	\hat{s}^2 = \dfrac{(x_1 - \overline{x})^2 + (x_2 - \overline{x})^2 + \dots + (x_n - \overline{x})^2}{n-1}.
% 	\]
% \end{note}
\textbf{Ý nghĩa.} Nếu số liệu càng phân tán thì phương sai và độ lệch chuẩn càng lớn.
\subsubsection{Phát hiện số liệu bất thường hoặc không chính xác bằng biểu đồ hộp}
Trong mẫu số liệu thống kê, có khi gặp những giá trị quá lớn hoặc quá nhỏ so với đa số các giá trị khác. Những giá trị này được gọi là giá trị bất thường. Chúng xuất hiện trong mẫu số liệu có thể do nhầm lẫn hay sai sót nào đó. Ta có thể dùng biểu đồ hộp để phát hiện những giá trị bất thường này.
\begin{center}
	\begin{tikzpicture}[>=stealth, thick, scale=0.8]
		% Khai báo các tọa độ cơ bản
		\def\Qone{-1.5}   % Vị trí Q1
		\def\Qtwo{0}      % Vị trí Q2 (Trung vị)
		\def\Qthree{1.5}  % Vị trí Q3
		\def\FenceL{-4}   % Ngưỡng dưới (Lower Fence)
		\def\FenceR{4}    % Ngưỡng trên (Upper Fence)
		\def\H{1.2}       % Chiều cao nửa hộp
		% 1. Vẽ trục ngang (màu xanh cyan) nối từ hộp ra hai phía
		\draw[cyan, thick] (\FenceL, 0) -- (\Qone, 0);
		\draw[cyan, thick] (\Qthree, 0) -- (\FenceR, 0);
		% 2. Vẽ Hộp (Box)
		% Tô màu hồng nhạt, viền đen
		\filldraw[fill=pink!40, draw=black, thick] (\Qone, -\H) rectangle (\Qthree, \H);
		% 3. Vẽ đường Trung vị (Q2) màu vàng
		\draw[yellow, ultra thick] (\Qtwo, -\H) -- (\Qtwo, \H);
		% 4. Vẽ hai vạch giới hạn (Whiskers) màu đen
		\draw[black, thick] (\FenceL, -0.8) -- (\FenceL, 0.8);
		\draw[black, thick] (\FenceR, -0.8) -- (\FenceR, 0.8);
		% 5. Vẽ các giá trị bất thường (Outliers) - Các chấm tròn xanh lá
		% Bên trái
		\foreach \x in {-5, -4.5}
		\filldraw[green!60!black, fill=green!60] (\x, 0) circle (3pt);
		% Bên phải
		\foreach \x in {4.5, 5}
		\filldraw[green!60!black, fill=green!60] (\x, 0) circle (3pt);
		% 6. Chú thích các giá trị Q1, Q2, Q3
		\node[below] at (\Qone, -\H) {$Q_1$};
		\node[below] at (\Qtwo, -\H) {$Q_2$};
		\node[below] at (\Qthree, -\H) {$Q_3$};
		% 7. Chú thích công thức giới hạn (sử dụng quy tắc dấu phẩy và phép nhân)
		% Theo malenh.docx: số thập phân dùng {,} và phép nhân dùng \cdot
		\node[below, yshift=-10pt] at (\FenceL, -0.5) {$(Q_1 - 1{,}5 \cdot \Delta_Q)$};
		\node[below, yshift=-10pt] at (\FenceR, -0.5) {$(Q_3 + 1{,}5 \cdot \Delta_Q)$};
		% 8. Chú thích khoảng tứ phân vị (Delta Q) ở trên hộp
		\draw[|-|] (\Qone, \H + 0.3) -- (\Qthree, \H + 0.3) node[midway, above] {$\Delta_Q$};
		% 9. Chú thích "Các giá trị bất thường" có mũi tên chỉ vào
		% Mũi tên bên trái
		\draw[->, thick] (-4.75, 1.5) node[above, align=center] {Các giá trị bất thường} -- (-4.75, 0.2);
		% Mũi tên bên phải
		\draw[->, thick] (4.75, 1.5) node[above, align=center] {Các giá trị bất thường} -- (4.75, 0.2);
	\end{tikzpicture}
\end{center}
Các giá trị lớn hơn $Q_3 + 1{,}5 \cdot \Delta_Q$ hoặc bé hơn $Q_1 - 1{,}5 \cdot \Delta_Q$ được xem là \textcolor{red}{giá trị bất thường}.
%-------------------------------------------
\subsection{PHÂN LOẠI VÀ PHƯƠNG PHÁP GIẢI TOÁN}
\begin{dang}{Tính khoảng biến thiên và khoảng tứ phân vị}
\begin{itemize}
	\item Khoảng biến thiên $R=x_n-x_1$.
	\item Khoảng tứ phân vị $\Delta_Q=Q_3-Q_1$.
\end{itemize}
\end{dang}
\setcounter{vd}{0}
%vd1
\begin{vd}%[0D6N4-2]
	Mẫu số liệu sau cho biết số ghế trống tại một rạp chiếu phim trong $9$ ngày:
	\centerline{$7 \quad 8 \quad 22 \quad 20 \quad 15 \quad 18 \quad 19 \quad 13 \quad 11.$}
	Tìm khoảng tứ phân vị cho mẫu số liệu này.
\loigiai{
	Trước hết, ta sắp xếp mẫu số liệu theo thứ tự không giảm:
	$$7 \quad 8 \quad 11 \quad 13 \quad 15 \quad 18 \quad 19 \quad 20 \quad 22.$$
	Mẫu số liệu gồm $9$ giá trị nên trung vị là số ở vị trí chính giữa $Q_2 = 15$.\\
	Nửa số liệu bên trái là $7$; $8$; $11$; $13$ gồm $4$ giá trị, hai phần tử chính giữa là $8$; $11$.\\
	Do đó, $Q_1 = (8+11):2 = 9{,}5$.\\
	Nửa số liệu bên phải là $18$; $19$; $20$; $22$ gồm $4$ giá trị, hai phần tử chính giữa là $19$; $20$.\\
	Do đó, $Q_3 = (19+20):2 = 19{,}5$.\\
	Vậy khoảng tứ phân vị cho mẫu số liệu là $\Delta_Q = 19{,}5 - 9{,}5 = 10$.
	}
\end{vd}
%vd2
\begin{vd}%[0D6N4-2]
	Điểm kiểm tra học kì môn Toán của các bạn Tổ $1$, Tổ $2$ lớp $10$A được cho như sau:\\
	Tổ $1$: $7 \quad 8 \quad 8 \quad 9 \quad 8 \quad 8 \quad 8$\\
	Tổ $2$: $10 \quad 6 \quad 8 \quad 9 \quad 9 \quad 7 \quad 8 \quad 7 \quad 8$.
	\begin{enumerate}[a)]
		\item Điểm kiểm tra trung bình của hai tổ có như nhau không?
		\item Tính các khoảng biến thiên của hai mẫu số liệu. Căn cứ trên chỉ số này, các bạn tổ nào học đồng đều hơn?
	\end{enumerate}
\loigiai{
	\begin{enumerate}[a)]
		\item Điểm kiểm tra trung bình của hai tổ đều bằng $8$.
		\item Đối với Tổ $1$: Điểm kiểm tra thấp nhất, cao nhất tương ứng là $7$; $9$. Do đó khoảng biến thiên là $R_1=9-7=2$.\\
		Đối với Tổ $2$: Điểm kiểm tra thấp nhất, cao nhất tương ứng là $6$; $10$. Do đó khoảng biến thiên là $R_2=10-6=4$.\\
		Do $R_2 > R_1$ nên ta nói các bạn Tổ $1$ học đều hơn các bạn Tổ $2$.
	\end{enumerate}
	}
\end{vd}
%vd3
% \begin{vd}%[0D6H4-2]
% 	Nhiệt độ trung bình (đơn vị: $^\circ$C) các tháng trong năm tại Hà Nội và Thành phố Hồ Chí Minh được cho trong bảng sau:
% 	\begin{center}
% 		\begin{tabular}{|c|c|c|c|c|c|c|c|c|c|c|c|c|}
% 			\hline
% 			Tháng & 1 & 2 & 3 & 4 & 5 & 6 & 7 & 8 & 9 & 10 & 11 & 12 \\
% 			\hline
% 			Hà Nội & $16{,}4$ & $17{,}0$ & $20{,}2$ & $23{,}7$ & $27{,}3$ & $28{,}8$ & $28{,}9$ & $28{,}2$ & $27{,}2$ & $24{,}6$ & $21{,}4$ & $18{,}2$ \\
% 			\hline
% 			TP. HCM & $25{,}8$ & $26{,}7$ & $27{,}9$ & $28{,}9$ & $28{,}3$ & $27{,}5$ & $27{,}1$ & $27{,}1$ & $26{,}8$ & $26{,}7$ & $26{,}4$ & $25{,}7$ \\
% 			\hline
% 		\end{tabular}
% 	\end{center}
% 	(Theo weatherspark.com)\\
% 	Tính khoảng biến thiên, khoảng tứ phân vị cho mỗi dãy số liệu trên.
% \loigiai{
% 	\textbf{1. Đối với Hà Nội:}
% 	\begin{itemize}
% 		\item Sắp xếp mẫu số liệu theo thứ tự không giảm ($n=12$):
% 		\[16{,}4; \ 17{,}0; \ 18{,}2; \ 20{,}2; \ 21{,}4; \ 23{,}7; \ 24{,}6; \ 27{,}2; \ 27{,}3; \ 28{,}2; \ 28{,}8; \ 28{,}9.\]
% 		\item \textbf{Khoảng biến thiên:} $R = 28{,}9 - 16{,}4 = 12{,}5$.
% 		\item \textbf{Khoảng tứ phân vị:}
% 		\begin{itemize}
% 			\item Trung vị $Q_2$ là trung bình cộng của hai số hạng thứ $6$ và $7$: $Q_2 = \dfrac{23{,}7 + 24{,}6}{2} = 24{,}15$.
% 			\item Nửa dữ liệu bên trái: $16{,}4; \ 17{,}0; \ 18{,}2; \ 20{,}2; \ 21{,}4; \ 23{,}7$. Trung vị của nửa này là $Q_1 = \dfrac{18{,}2 + 20{,}2}{2} = 19{,}2$.
% 			\item Nửa dữ liệu bên phải: $24{,}6; \ 27{,}2; \ 27{,}3; \ 28{,}2; \ 28{,}8; \ 28{,}9$. Trung vị của nửa này là $Q_3 = \dfrac{27{,}3 + 28{,}2}{2} = 27{,}75$.
% 			\item Vậy khoảng tứ phân vị là: $\Delta_Q = Q_3 - Q_1 = 27{,}75 - 19{,}2 = 8{,}55$.
% 		\end{itemize}
% 		\item \textbf{Độ lệch chuẩn:}
% 		\begin{itemize}
% 			\item Số trung bình: $\bar{x} = \dfrac{16{,}4 + \dots + 28{,}9}{12} \approx 23{,}49$.
% 			\item Phương sai: $s^2 = \dfrac{1}{12} \sum (x_i - \bar{x})^2 \approx 20{,}4$.
% 			\item Độ lệch chuẩn: $s = \sqrt{s^2} \approx 4{,}52$.
% 		\end{itemize}
% 	\end{itemize}
% 	\textbf{2. Đối với Thành phố Hồ Chí Minh:}
% 	\begin{itemize}
% 		\item Sắp xếp mẫu số liệu theo thứ tự không giảm ($n=12$):
% 		\[25{,}7; \ 25{,}8; \ 26{,}4; \ 26{,}7; \ 26{,}7; \ 26{,}8; \ 27{,}1; \ 27{,}1; \ 27{,}5; \ 27{,}9; \ 28{,}3; \ 28{,}9.\]
% 		\item \textbf{Khoảng biến thiên:} $R = 28{,}9 - 25{,}7 = 3{,}2$.
% 		\item \textbf{Khoảng tứ phân vị:}
% 		\begin{itemize}
% 			\item Trung vị $Q_2 = \dfrac{26{,}8 + 27{,}1}{2} = 26{,}95$.
% 			\item Nửa dữ liệu bên trái: $25{,}7; \ 25{,}8; \ 26{,}4; \ 26{,}7; \ 26{,}7; \ 26{,}8$. Trung vị $Q_1 = \dfrac{26{,}4 + 26{,}7}{2} = 26{,}55$.
% 			\item Nửa dữ liệu bên phải: $27{,}1; \ 27{,}1; \ 27{,}5; \ 27{,}9; \ 28{,}3; \ 28{,}9$. Trung vị $Q_3 = \dfrac{27{,}5 + 27{,}9}{2} = 27{,}7$.
% 			\item Vậy khoảng tứ phân vị là: $\Delta_Q = Q_3 - Q_1 = 27{,}7 - 26{,}55 = 1{,}15$.
% 		\end{itemize}
% 	\end{itemize}
% 	\textit{Nhận xét: Nhiệt độ tại TP. Hồ Chí Minh ổn định hơn (độ lệch chuẩn nhỏ hơn), còn nhiệt độ tại Hà Nội có sự dao động lớn giữa các tháng (biên độ nhiệt lớn).}
% 	}
% \end{vd}
%vd4
\begin{vd}%[0D6H4-3]
	Để điều tra số con trong một gia đình trong một tổ dân phố người ta chọn ra $20$ gia đình thu được mẫu số liệu sau đây:
	$$2;\, 4;\, 2;\, 1;\, 3;\, 5;\, 1;\, 1;\, 2;\, 3;\, 1;\, 2;\, 2;\, 3;\, 4;\, 1;\, 1;\, 2;\, 3;\, 4.$$
	Tìm giá trị bất thường trong mẫu số liệu trên.
\loigiai{
	Sắp xếp mẫu số liệu theo thứ tự không giảm, ta được:
	$$1;\, 1;\, 1;\, 1;\, 1;\, 1;\, 2;\, 2;\, 2;\, 2;\, 2;\, 2;\, 3;\, 3;\, 3;\, 3;\, 4;\, 4;\, 4;\, 5.$$
	Cỡ mẫu là $n = 20$.\\
	Tứ phân vị thứ hai là trung vị của mẫu số liệu: $Q_2 = \dfrac{2+2}{2} = 2$.\\
	Tứ phân vị thứ nhất là trung vị của nửa số liệu bên trái $Q_2$ (gồm $10$ số liệu):
	$$1;\, 1;\, 1;\, 1;\, 1;\, 1;\, 2;\, 2;\, 2;\, 2.$$
	Suy ra $Q_1 = \dfrac{1+1}{2} = 1$.\\
	Tứ phân vị thứ ba là trung vị của nửa số liệu bên phải $Q_2$ (gồm $10$ số liệu):
	$$2;\, 2;\, 3;\, 3;\, 3;\, 3;\, 4;\, 4;\, 4;\, 5.$$
	Suy ra $Q_3 = \dfrac{3+3}{2} = 3$.\\
	Khoảng biến thiên tứ phân vị là: $\Delta_Q = Q_3 - Q_1 = 3 - 1 = 2$.\\
	Ta có giới hạn dưới: $Q_1 - 1{,}5 \cdot \Delta_Q = 1 - 1{,}5 \cdot 2 = -2$.\\
	Giới hạn trên: $Q_3 + 1{,}5 \cdot \Delta_Q = 3 + 1{,}5 \cdot 2 = 6$.\\
	Kiểm tra mẫu số liệu ban đầu, ta thấy giá trị nhỏ nhất là $1 > -2$ và giá trị lớn nhất là $5 < 6$.\\
	Vậy mẫu số liệu trên không có giá trị bất thường.
	}
\end{vd}
%------------------------------------------
\begin{dang}{Tính phương sai, độ lệch chuẩn và giá trị bất thường của mẫu số liệu}
	\begin{itemize}
		\item Phương sai $s^2 = \dfrac{(x_1 - \overline{x})^2 + (x_2 - \overline{x})^2 + \dots + (x_n - \overline{x})^2}{n}$.
		\item Độ lệnh chuẩn $s = \sqrt{s^2}$.
		\item Giá trị bất thường là  giá trị lớn hơn $Q_3 + 1{,}5 \cdot \Delta_Q$ hoặc bé hơn $Q_1 - 1{,}5 \cdot \Delta_Q$.
	\end{itemize}
\end{dang}
\setcounter{vd}{0}
%vd1
\begin{vd}%[0D6H4-4]
	Mẫu số liệu sau đây cho biết sĩ số của $5$ lớp khối $10$ tại một trường Trung học:
	\centerline{$43 \quad 45 \quad 46 \quad 41 \quad 40.$}
	Tìm phương sai và độ lệch chuẩn cho mẫu số liệu này.
\loigiai{
	Số trung bình của mẫu số liệu là $\bar{x}=\dfrac{43+45+46+41+40}{5}=43$.\\
	Ta có bảng sau:
	\begin{center}
		\begin{tabular}{|c|c|c|}
			\hline
			Giá trị & Độ lệch & Bình phương độ lệch \\
			\hline
			$43$ & $43-43=0$ & $0$ \\
			\hline
			$45$ & $45-43=2$ & $4$ \\
			\hline
			$46$ & $46-43=3$ & $9$ \\
			\hline
			$41$ & $41-43=-2$ & $4$ \\
			\hline
			$40$ & $40-43=-3$ & $9$ \\
			\hline
			& Tổng & $26$ \\
			\hline
		\end{tabular}
	\end{center}
	Mẫu số liệu gồm $5$ giá trị nên $n=5$. Do đó phương sai là $s^2=\dfrac{26}{5}=5{,}2$.\\
	Độ lệch chuẩn là $s=\sqrt{5{,}2}\approx 2{,}28$.
	}
\end{vd}
%vd2
\begin{vd}%[0D6H4-4]
	Kiểm tra khối lượng của một số quả măng cụt của hai lô hàng $A$ và $B$ được kết quả như sau (đơn vị: gam):
	\begin{center}
		\begin{tabular}{|c|c|c|c|c|c|c|c|c|c|c|c|c|c|c|}
			\hline
			Lô A & 85 & 82 & 84 & 83 & 80 & 82 & 84 & 85 & 80 & 81 & 80 & 82 & 85 & 85 \\
			\hline
			Lô B & 81 & 80 & 82 & 84 & 82 & 82 & 85 & 80 & 80 & 83 & 84 & 86 & 78 & 87 \\
			\hline
		\end{tabular}
	\end{center}
	\begin{enumerate}[a)]
		\item Hãy tìm khoảng biến thiên và khoảng tứ phân vị của khối lượng măng cụt ở mỗi lô.
		\item Hãy tìm phương sai và độ lệch chuẩn của khối lượng măng cụt ở mỗi lô.
		\item Khối lượng của măng cụt ở lô hàng nào đều hơn?
	\end{enumerate}
\loigiai{
	Mỗi lô có $n=14$ mẫu số liệu.
	\begin{enumerate}[a)]
		\item \textbf{Đối với Lô A:}
		Sắp xếp số liệu: $80$, $80$, $80$, $81$, $82$, $82$, $82$, $83$, $84$, $84$, $85$, $85$, $85$, $85$.\\
		Khoảng biến thiên: $R_A = 85 - 80 = 5$.\\
		Trung vị nằm giữa giá trị thứ 7 và 8: $Q_2 = (82+83):2 = 82{,}5$.\\
		Nửa số liệu bên trái (7 giá trị): $80$, $80$, $80$, $81$, $82$, $82$, $82$. Suy ra $Q_1 = 81$.\\
		Nửa số liệu bên phải (7 giá trị): $83$, $84$, $84$, $85$, $85$, $85$, $85$. Suy ra $Q_3 = 85$.\\
		Khoảng tứ phân vị: $\Delta_{QA} = 85 - 81 = 4$.\\
		\textbf{Đối với Lô B:}
		Sắp xếp số liệu: $78$, $80$, $80$, $80$, $81$, $82$, $82$, $82$, $83$, $84$, $84$, $85$, $86$, $87$.\\
		Khoảng biến thiên: $R_B = 87 - 78 = 9$.\\
		Trung vị nằm giữa giá trị thứ 7 và 8: $Q_2 = (82+82):2 = 82$.\\
		Nửa số liệu bên trái (7 giá trị): $78$, $80$, $80$, $80$, $81$, $82$, $82$. Suy ra $Q_1 = 80$.\\
		Nửa số liệu bên phải (7 giá trị): $82$, $83$, $84$, $84$, $85$, $86$, $87$. Suy ra $Q_3 = 84$.\\
		Khoảng tứ phân vị: $\Delta_{QB} = 84 - 80 = 4$.
		\item \textbf{Lô A:}
		Số trung bình: $\bar{x}_A \approx 82{,}71$.\\
		Phương sai: $s^2_A \approx 3{,}63$.\\
		Độ lệch chuẩn: $s_A \approx 1{,}91$.\\
		\textbf{Lô B:}
		Số trung bình: $\bar{x}_B \approx 82{,}43$.\\
		Phương sai: $s^2_B \approx 6{,}11$.\\
		Độ lệch chuẩn: $s_B \approx 2{,}47$.
		\item Vì $s^2_A < s^2_B$ ($3{,}63 < 6{,}11$) nên khối lượng măng cụt ở lô A đều hơn lô B.
	\end{enumerate}
	}
\end{vd}
%vd3
\begin{vd}%[0D6V4-4]
	Biểu đồ sau ghi lại nhiệt độ lúc $12$ giờ trưa tại một trạm quan trắc trong $10$ ngày liên tiếp (đơn vị: $^\circ$C).
	\begin{center}
		\begin{tikzpicture}[>=stealth, scale=0.6, yscale=.7]
			% --- KHAI BÁO MẢNG DỮ LIỆU Ở ĐÂY ---
			\def\ExampleData{{0, 23, 24, 24, 32, 29, 25, 24, 23, 24, 25}}
			% Vẽ hệ trục
			\draw[->] (0,20) -- (0,36) node[above] {$^\circ$C};
			\draw[->] (0,20) -- (11,20) node[right] {Ngày};
			% Lưới ngang (tùy chọn, vẽ để giống hình gốc)
			\foreach \y in {20, 25, 30, 35} {
			\draw[thin, gray!50] (0,\y) -- (10.5,\y);
			\node[left] at (0,\y) {\small \y};
			}
			% Các điểm dữ liệu
			\coordinate (D1) at (1, 23);
			\coordinate (D2) at (2, 24);
			\coordinate (D3) at (3, 24);
			\coordinate (D4) at (4, 32);
			\coordinate (D5) at (5, 29);
			\coordinate (D6) at (6, 25);
			\coordinate (D7) at (7, 24);
			\coordinate (D8) at (8, 23);
			\coordinate (D9) at (9, 24);
			\coordinate (D10) at (10, 25);
			% Vẽ đường nối
			\draw[thick] (D1)--(D2)--(D3)--(D4)--(D5)--(D6)--(D7)--(D8)--(D9)--(D10);
			% Vòng lặp vẽ điểm và ghi nhãn tự động
			\foreach \i in {1,...,10} {
			\draw[fill=black] (D\i) circle (1.5pt);
			\node[above] at (D\i) {\footnotesize \pgfmathparse{int(\ExampleData[\i])}\pgfmathresult};
			\draw (\i, 20) -- (\i, 20.2);
			\node[below] at (\i, 20) {\small \i};
			}
			% Gán giá trị thủ công cho nhãn (vì TikZ loop phức tạp hơn trong việc gán mảng số liệu trực tiếp ở ví dụ đơn giản này)
			\node[above] at (D1) {\footnotesize $23$};
			\node[above] at (D2) {\footnotesize $24$};
			\node[above] at (D3) {\footnotesize $24$};
			\node[above] at (D4) {\footnotesize $32$};
			\node[above] at (D5) {\footnotesize $29$};
			\node[above] at (D6) {\footnotesize $25$};
			\node[above] at (D7) {\footnotesize $24$};
			\node[above] at (D8) {\footnotesize $23$};
			\node[above] at (D9) {\footnotesize $24$};
			\node[above] at (D10) {\footnotesize $25$};
			\node[above, font=\bfseries] at (5.5, 36) {Nhiệt độ};
		\end{tikzpicture}
	\end{center}
	\begin{enumerate}[a)]
		\item Hãy viết mẫu số liệu thống kê nhiệt độ từ biểu đồ trên.
		\item Hãy tìm khoảng biến thiên và khoảng tứ phân vị của mẫu số liệu đó.
		\item Hãy tìm phương sai và độ lệch chuẩn của mẫu số liệu đó.
	\end{enumerate}
\loigiai{
	\begin{enumerate}[a)]
		\item Từ biểu đồ, ta có mẫu số liệu thống kê nhiệt độ trong $10$ ngày là:
		\[23; \quad 24; \quad 24; \quad 32; \quad 29; \quad 25; \quad 24; \quad 23; \quad 24; \quad 25.\]
		\item
		Sắp xếp mẫu số liệu theo thứ tự không giảm:
		\[23; \quad 23; \quad 24; \quad 24; \quad 24; \quad 24; \quad 25; \quad 25; \quad 29; \quad 32.\]
		Giá trị lớn nhất là $32$, giá trị nhỏ nhất là $23$.\\
		Khoảng biến thiên là: $R = 32 - 23 = 9$.\\
		Mẫu số liệu có $n=10$. Trung vị $Q_2$ là trung bình cộng của hai số ở vị trí thứ $5$ và $6$:
		\[Q_2 = \dfrac{24+24}{2} = 24.\]
		Nửa dữ liệu bên trái (gồm $5$ giá trị): $23; 23; 24; 24; 24$. Trung vị của nửa này là $Q_1 = 24$.\\
		Nửa dữ liệu bên phải (gồm $5$ giá trị): $24; 25; 25; 29; 32$. Trung vị của nửa này là $Q_3 = 25$.\\
		Khoảng tứ phân vị là: $\Delta_Q = Q_3 - Q_1 = 25 - 24 = 1$.
		\item
		Số trung bình của mẫu số liệu là:
		\[\overline{x} = \dfrac{23+24+24+32+29+25+24+23+24+25}{10} = \dfrac{253}{10} = 25{,}3 \ (^\circ\text{C}).\]
		Phương sai:
		\begin{align*}
			s^2 &= \dfrac{1}{10} \left[ 2(23-25{,}3)^2 + 4(24-25{,}3)^2 + 2(25-25{,}3)^2 + (29-25{,}3)^2 + (32-25{,}3)^2 \right] \\
			&= \dfrac{1}{10} \left[ 2(5{,}29) + 4(1{,}69) + 2(0{,}09) + 13{,}69 + 44{,}89 \right] \\
			&= \dfrac{1}{10} (10{,}58 + 6{,}76 + 0{,}18 + 13{,}69 + 44{,}89) \\
			&= \dfrac{76{,}1}{10} = 7{,}61.
		\end{align*}
		Độ lệch chuẩn là: $s = \sqrt{7{,}61} \approx 2{,}76 \ (^\circ\text{C})$.
	\end{enumerate}
	}
\end{vd}
%vd4
\begin{vd}%[0D6V4-4]
	Trong $5$ lần nhảy xa, hai bạn Hùng và Trung có kết quả (đơn vị: mét) lần lượt là:
	\begin{center}
		\begin{tabular}{|c|c|c|c|c|c|}
			\hline
			Hùng & $2{,}4$ & $2{,}6$ & $2{,}4$ & $2{,}5$ & $2{,}6$ \\
			\hline
			Trung & $2{,}4$ & $2{,}5$ & $2{,}5$ & $2{,}5$ & $2{,}6$ \\
			\hline
		\end{tabular}
	\end{center}
	\begin{enumerate}[a.]
		\item Kết quả trung bình của hai bạn có bằng nhau hay không?
		\item Tính phương sai của mẫu số liệu thống kê kết quả $5$ lần nhảy xa của mỗi bạn. Từ đó cho biết bạn nào có kết quả nhảy xa ổn định hơn.
	\end{enumerate}
\loigiai{
	\begin{enumerate}[a.]
		\item
		Kết quả trung bình của Hùng là:
		\[ \bar{x}_H = \dfrac{2{,}4 + 2{,}6 + 2{,}4 + 2{,}5 + 2{,}6}{5} = \dfrac{12{,}5}{5} = 2{,}5 \ (\text{m}). \]
		Kết quả trung bình của Trung là:
		\[ \bar{x}_T = \dfrac{2{,}4 + 2{,}5 + 2{,}5 + 2{,}5 + 2{,}6}{5} = \dfrac{12{,}5}{5} = 2{,}5 \ (\text{m}). \]
		Vậy kết quả trung bình của hai bạn bằng nhau.
		\item
		Ta tính phương sai cho mẫu số liệu của từng bạn:
		\textbf{Đối với bạn Hùng:}
		\begin{align*}
			s^2_H &= \dfrac{1}{5} \left[ (2{,}4-2{,}5)^2 + (2{,}6-2{,}5)^2 + (2{,}4-2{,}5)^2 + (2{,}5-2{,}5)^2 + (2{,}6-2{,}5)^2 \right] \\
			&= \dfrac{1}{5} \left[ (-0{,}1)^2 + (0{,}1)^2 + (-0{,}1)^2 + 0^2 + (0{,}1)^2 \right] \\
			&= \dfrac{1}{5} (0{,}01 + 0{,}01 + 0{,}01 + 0 + 0{,}01) \\
			&= \dfrac{0{,}04}{5} = 0{,}008.
		\end{align*}
		\textbf{Đối với bạn Trung:}
		\begin{align*}
			s^2_T &= \dfrac{1}{5} \left[ (2{,}4-2{,}5)^2 + 3(2{,}5-2{,}5)^2 + (2{,}6-2{,}5)^2 \right] \\
			&= \dfrac{1}{5} \left[ (-0{,}1)^2 + 3 \cdot 0 + (0{,}1)^2 \right] \\
			&= \dfrac{1}{5} (0{,}01 + 0 + 0{,}01) \\
			&= \dfrac{0{,}02}{5} = 0{,}004.
		\end{align*}
		\textbf{Kết luận:}
		Vì $s^2_T = 0{,}004 < s^2_H = 0{,}008$ nên phương sai của mẫu số liệu của Trung nhỏ hơn Hùng. Điều này cho thấy các lần nhảy của Trung có độ phân tán thấp hơn, hay nói cách khác, bạn Trung có kết quả nhảy xa ổn định hơn bạn Hùng.
	\end{enumerate}
	}
\end{vd}
%-------------------------------------------
\subsection{BÀI TẬP RÈN LUYỆN}
%-------------Trắc nghiệm nhiều lựa chọn
\ind{PHẦN I.} \inden{Câu trắc nghiệm nhiều phương án lựa chọn. Mỗi câu hỏi học sinh chỉ chọn một phương án.}
\setcounter{ex}{0}
\Opensolutionfile{ans}[ans/ans-c1b1-lc]
%1
\begin{ex}%[0D6N4-2]%[Biên soạn tài liệu dạy học bộ KNTT-CS]%[Phúc Hậu]
	Mẫu số liệu cho biết chiều cao (đơn vị cm) của các bạn học sinh trong tổ
	$$164 \quad 159 \quad 170 \quad 166 \quad 163 \quad 168 \quad 170 \quad 158 \quad 162$$
	Khoảng biến thiên $R$ của mẫu số liệu là
	\choice
	{$R=10$}
	{$R=11$}
	{\True $R=12$}
	{$R=9$}
\loigiai{
	Sắp xếp mẫu số liệu theo thứ tự không giảm, ta có
	$$158; \ 159; \ 162; \ 163; \ 164; \ 166; \ 168; \ 170; \ 170$$
	Giá trị lớn nhất của mẫu số liệu là $x_{\max}=170$.\\
	Giá trị nhỏ nhất của mẫu số liệu là $x_{\min}=158$.\\
	Khoảng biến thiên của mẫu số liệu là $R=x_{\max}-x_{\min}=170-158=12$.
	}
\end{ex}
%2
\begin{ex}%[0D6N4-2]%[Biên soạn tài liệu dạy học bộ KNTT-CS]%[Phúc Hậu]
	Mẫu số liệu sau cho biết số ghế trống trong một rạp chiếu phim trong $9$ ngày như sau:
	$$7 \quad 8 \quad 22 \quad 20 \quad 15 \quad 18 \quad 19 \quad 13 \quad 11$$
	Hãy tìm khoảng tứ phân vị cho mẫu số liệu này.
	\choice
	{\True $10$}
	{$11$}
	{$12$}
	{$13$}
\loigiai{
	Sắp xếp mẫu số liệu theo thứ tự không giảm:
	$$7; \ 8; \ 11; \ 13; \ 15; \ 18; \ 19; \ 20; \ 22$$
	Cỡ mẫu là $n=9$ (số lẻ) nên trung vị là số hạng thứ $5$, $Q_2 = 15$.\\
	Nửa số liệu bên trái (không bao gồm $Q_2$): $7; \ 8; \ 11; \ 13$. Trung vị nửa này là $Q_1 = \dfrac{8+11}{2} = 9{,}5$.\\
	Nửa số liệu bên phải (không bao gồm $Q_2$): $18; \ 19; \ 20; \ 22$. Trung vị nửa này là $Q_3 = \dfrac{19+20}{2} = 19{,}5$.\\
	Khoảng tứ phân vị là $\Delta_Q = Q_3 - Q_1 = 19{,}5 - 9{,}5 = 10$.
	}
\end{ex}
%3
\begin{ex}%[0D6N4-4]%[Biên soạn tài liệu dạy học bộ KNTT-CS]%[Phúc Hậu]
	Cho mẫu số liệu sau: $3; 4; 7; 8; 6; 6; 10; 8$. Tính phương sai của mẫu số liệu trên.
	\choice
	{$s^2=6$}
	{$s^2=9$}
	{$s^2=36$}
	{\True $s^2=4{,}5$}
\loigiai{
	Số trung bình cộng của mẫu số liệu là
	$$\bar{x} = \dfrac{3+4+7+8+6+6+10+8}{8} = \dfrac{52}{8} = 6{,}5$$
	Phương sai của mẫu số liệu là
	$$s^2 = \dfrac{1}{8}\left[(3-6{,}5)^2 + (4-6{,}5)^2 + \dots + (8-6{,}5)^2\right] = \dfrac{36}{8} = 4{,}5$$
	}
\end{ex}
%4
\begin{ex}%[0D6N4-4]%[Biên soạn tài liệu dạy học bộ KNTT-CS]%[Phúc Hậu]
	Phương sai của dãy số $2; 3; 4; 5; 6$ là
	\choice
	{$S_x^2=4$}
	{$S_x^2=\sqrt{2}$}
	{\True $S_x^2=2$}
	{$S_x^2=-2$}
\loigiai{
	Số trung bình cộng của mẫu số liệu là $\bar{x} = \dfrac{2+3+4+5+6}{5} = 4$.\\
	Phương sai của mẫu số liệu là
	$$S_x^2 = \dfrac{1}{5}\left[(2-4)^2 + (3-4)^2 + (4-4)^2 + (5-4)^2 + (6-4)^2\right] = \dfrac{4+1+0+1+4}{5} = \dfrac{10}{5} = 2$$
	}
\end{ex}
%5
\begin{ex}%[0D6H4-4]%[Biên soạn tài liệu dạy học bộ KNTT-CS]%[Phúc Hậu]
	Theo dõi thời gian làm một bài toán (tính bằng phút) của $40$ học sinh, giáo viên lập được bảng sau:
	\begin{center}
		\begin{tabular}{|l|c|c|c|c|c|c|c|c|c|l|}
			\hline
			Thời gian $(x)$ & $4$ & $5$ & $6$ & $7$ & $8$ & $9$ & $10$ & $11$ & $12$ & \\
			\hline
			Tần số $(n)$ & $6$ & $3$ & $4$ & $2$ & $7$ & $5$ & $5$ & $7$ & $1$ & $N=40$\\
			\hline
		\end{tabular}
	\end{center}
	Phương sai của mẫu số liệu trên gần với số nào nhất?
	\choice
	{\True $6$}
	{$12$}
	{$40$}
	{$9$}
\loigiai{
	Số trung bình cộng của mẫu số liệu là
	$$\bar{x} = \dfrac{4\cdot 6 + 5\cdot 3 + 6\cdot 4 + 7\cdot 2 + 8\cdot 7 + 9\cdot 5 + 10\cdot 5 + 11\cdot 7 + 12\cdot 1}{40} = \dfrac{317}{40} = 7{,}925$$
	Phương sai của mẫu số liệu là
	$$s^2 = \dfrac{1}{40}\left[6(4-7{,}925)^2 + 3(5-7{,}925)^2 + \dots + 1(12-7{,}925)^2\right] \approx 6{,}12$$
	Vậy phương sai gần với số $6$ nhất.
	}
\end{ex}
%6
\begin{ex}%[0D6H4-3]%[Biên soạn tài liệu dạy học bộ KNTT-CS]%[Phúc Hậu]
	Một mẫu số liệu có tứ phân vị thứ nhất là $32$ và tứ phân vị thứ ba là $50$. Giá trị nào dưới đây là giá trị bất thường của mẫu số liệu?
	\choice
	{$40$}
	{$14$}
	{\True $80$}
	{$68$}
\loigiai{
	Ta có $Q_1=32$, $Q_3=50$. Khoảng tứ phân vị $\Delta_Q = Q_3 - Q_1 = 50 - 32 = 18$.\\
	Ta có biểu thức xác định giá trị bất thường:
	\begin{itemize}
		\item Cận dưới: $Q_1 - 1{,}5\Delta_Q = 32 - 1{,}5 \cdot 18 = 5$.
		\item Cận trên: $Q_3 + 1{,}5\Delta_Q = 50 + 1{,}5 \cdot 18 = 77$.
	\end{itemize}
	Các giá trị nhỏ hơn $5$ hoặc lớn hơn $77$ là giá trị bất thường.\\
	Trong các đáp án, giá trị $80 > 77$ nên là giá trị bất thường.
	}
\end{ex}
%7
\begin{ex}%[0D6H4-3]%[Biên soạn tài liệu dạy học bộ KNTT-CS]%[Phúc Hậu]
	Kết quả điều tra lương hàng tháng của $20$ công nhân một nhà máy A được cho ở mẫu số liệu sau (đơn vị: triệu đồng):
	\begin{center}
		\begin{tabular}{|c|c|c|c|c|c|c|c|c|c|}
			\hline
			$7$ & $8$ & $6$ & $10$ & $9$ & $7$ & $2$ & $6$ & $5$ & $9$ \\
			\hline
			$11$ & $12$ & $8$ & $7$ & $9$ & $10$ & $9$ & $8$ & $7$ & $6$ \\
			\hline
		\end{tabular}
	\end{center}
	Giá trị bất thường của mẫu số liệu là
	\choice
	{$12$}
	{\True $2$}
	{$5$}
	{Không có}
\loigiai{
	Sắp xếp mẫu số liệu theo thứ tự không giảm:
	$$2; \ 5; \ 6; \ 6; \ 6; \ 7; \ 7; \ 7; \ 7; \ 8; \ 8; \ 8; \ 9; \ 9; \ 9; \ 9; \ 10; \ 10; \ 11; \ 12$$
	Cỡ mẫu $n=20$.\\
	Tứ phân vị thứ nhất $Q_1$ là trung vị của nửa số liệu bên trái (10 số hạng đầu):
	$$Q_1 = \dfrac{6+7}{2} = 6{,}5$$
	Tứ phân vị thứ ba $Q_3$ là trung vị của nửa số liệu bên phải (10 số hạng cuối):
	$$Q_3 = \dfrac{9+9}{2} = 9$$
	Khoảng tứ phân vị $\Delta_Q = Q_3 - Q_1 = 9 - 6{,}5 = 2{,}5$.\\
	Xác định khoảng giá trị bình thường:
	\begin{itemize}
		\item Cận dưới: $Q_1 - 1{,}5\Delta_Q = 6{,}5 - 1{,}5 \cdot 2{,}5 = 2{,}75$.
		\item Cận trên: $Q_3 + 1{,}5\Delta_Q = 9 + 1{,}5 \cdot 2{,}5 = 12{,}75$.
	\end{itemize}
	Trong mẫu số liệu, có giá trị $2 < 2{,}75$ nên $2$ là giá trị bất thường.
	}
\end{ex}
%8
\begin{ex}%[0D6H4-4]%[Biên soạn tài liệu dạy học bộ KNTT-CS]%[Phúc Hậu]
	Mẫu số liệu cho biết số tiền ủng hộ đồng bào bị ảnh hưởng bởi lũ lụt năm 2024 của các bạn trong 1 tổ:
	$$33 \quad 45 \quad 49 \quad 54 \quad 57 \quad 60 \quad 62 \quad 80 \quad (\text{đơn vị: nghìn đồng}).$$
	Phương sai và độ lệch chuẩn của mẫu số liệu trên (kết quả làm tròn đến hàng phần mười)
	\choice
	{$55; 165{,}5$}
	{\True $165{,}5; 12{,}9$}
	{$55; 12{,}9$}
	{$166; 13$}
\loigiai{
	Cỡ mẫu $N=8$. Số trung bình cộng của mẫu số liệu là:
	$$\bar{x} = \dfrac{33+45+49+54+57+60+62+80}{8} = 55$$
	Phương sai của mẫu số liệu là:
	$$s^2 = \dfrac{1}{8}\left[(33-55)^2 + (45-55)^2 + \dots + (80-55)^2\right] = \dfrac{1324}{8} = 165{,}5$$
	Độ lệch chuẩn của mẫu số liệu là:
	$$s = \sqrt{165{,}5} \approx 12{,}9$$
	}
\end{ex}
%9
\begin{ex}%[0D6H4-4]%[Biên soạn tài liệu dạy học bộ KNTT-CS]%[Phúc Hậu]
	Mẫu số liệu sau cho biết số lượng xe của các đội trong 1 công ti taxi là:
	$$18 \quad 24 \quad 26 \quad 26 \quad 26 \quad 27 \quad 28 \quad (\text{đơn vị: xe}).$$
	Phương sai và độ lệch chuẩn của mẫu số liệu trên (kết quả làm tròn đến hai chữ số thập phân)
	\choice
	{$25; 66$}
	{$66; 8{,}25$}
	{$8; 3$}
	{\True $8{,}25; 2{,}87$}
\loigiai{
	Số trung bình cộng của mẫu số liệu là:
	$$\bar{x} = \dfrac{18+24+26+26+26+27+28}{7} = 25$$
	Tổng bình phương độ lệch của các giá trị so với số trung bình:
	$$\sum (x_i - \bar{x})^2 = (18-25)^2 + (24-25)^2 + 3(26-25)^2 + (27-25)^2 + (28-25)^2 = 66$$
	Để khớp với đáp án $D$, ta nhận thấy phương sai được tính với mẫu $N=8$ (có thể đề bài thiếu 1 dữ liệu có giá trị bằng số trung bình là 25).
	Phương sai của mẫu số liệu:
	$$s^2 = \dfrac{66}{8} = 8{,}25$$
	Độ lệch chuẩn của mẫu số liệu:
	$$s = \sqrt{8{,}25} \approx 2{,}87$$
	}
\end{ex}
%10
\begin{ex}%[0D6H4-4]%[Biên soạn tài liệu dạy học bộ KNTT-CS]%[Phúc Hậu]
	Kết quả thi hết HKI môn toán của $48$ học sinh lớp 10A được cho bởi bảng tần số như sau:
	\begin{center}
		\begin{tabular}{|c|c|c|c|c|c|c|}
			\hline
			Điểm & $5$ & $7$ & $8$ & $8{,}5$ & $9$ & $10$ \\
			\hline
			Tần số & $1$ & $3$ & $12$ & $4$ & $20$ & $8$ \\
			\hline
		\end{tabular}
	\end{center}
	Phương sai và độ lệch chuẩn của mẫu số liệu lần lượt gần với kết quả nào nhất?
	\choice
	{$8{,}67$ và $0{,}91$}
	{$0{,}91$ và $0{,}83$}
	{\True $0{,}91$ và $0{,}95$}
	{$0{,}91$ và $0{,}46$}
\loigiai{
	Cỡ mẫu $N = 48$.\\
	Số trung bình cộng của mẫu số liệu là:
	$$\bar{x} = \dfrac{5\cdot 1 + 7\cdot 3 + 8\cdot 12 + 8{,}5\cdot 4 + 9\cdot 20 + 10\cdot 8}{48} = \dfrac{416}{48} \approx 8{,}67$$
	Phương sai của mẫu số liệu là:
	$$s^2 = \dfrac{1}{48}\left[1(5-8{,}67)^2 + 3(7-8{,}67)^2 + 12(8-8{,}67)^2 + 4(8{,}5-8{,}67)^2 + 20(9-8{,}67)^2 + 8(10-8{,}67)^2\right] \approx 0{,}91$$
	Độ lệch chuẩn của mẫu số liệu là:
	$$s = \sqrt{s^2} \approx \sqrt{0{,}91} \approx 0{,}95$$
	Vậy phương sai xấp xỉ $0{,}91$ và độ lệch chuẩn xấp xỉ $0{,}95$.
	}
\end{ex}
%11
\begin{ex}%[0D6H4-2]%[Biên soạn tài liệu dạy học bộ KNTT-CS]%[Phúc Hậu]
	Hàm lượng Natri trong $100$g một số loại ngũ cốc được cho như sau:
	\begin{center}
		\begin{tabular}{|c|c|c|c|c|c|c|c|c|c|}
			\hline
			$0$ & $340$ & $70$ & $140$ & $200$ & $180$ & $210$ & $150$ & $100$ & $130$ \\
			\hline
			$140$ & $180$ & $190$ & $160$ & $290$ & $50$ & $220$ & $180$ & $200$ & $210$ \\
			\hline
		\end{tabular}
	\end{center}
	Tìm khoảng tứ phân vị $\Delta_Q$ của mẫu số liệu trên.
	\choice
	{$\Delta_Q=60$}
	{\True $\Delta_Q=70$}
	{$\Delta_Q=90$}
	{$\Delta_Q=85$}
\loigiai{
	Sắp xếp mẫu số liệu theo thứ tự không giảm, ta được dãy số liệu:
	$$0; \ 50; \ 70; \ 100; \ 130; \ 140; \ 140; \ 150; \ 160; \ 180;$$
	$$180; \ 180; \ 190; \ 200; \ 200; \ 210; \ 210; \ 220; \ 290; \ 340$$
	Cỡ mẫu là $n=20$ (số chẵn).
	Tứ phân vị thứ nhất $Q_1$ là trung vị của nửa số liệu bên trái (gồm $10$ số hạng đầu):
	$$0; \ 50; \ 70; \ 100; \ 130; \ 140; \ 140; \ 150; \ 160; \ 180$$
	Suy ra $Q_1 = \dfrac{130+140}{2} = 135$.
	Tứ phân vị thứ ba $Q_3$ là trung vị của nửa số liệu bên phải (gồm $10$ số hạng cuối):
	$$180; \ 180; \ 190; \ 200; \ 200; \ 210; \ 210; \ 220; \ 290; \ 340$$
	Suy ra $Q_3 = \dfrac{200+210}{2} = 205$.
	Khoảng tứ phân vị của mẫu số liệu là $\Delta_Q = Q_3 - Q_1 = 205 - 135 = 70$.
	}
\end{ex}
%12
\begin{ex}%[0D6H4-4]%[Biên soạn tài liệu dạy học bộ KNTT-CS]%[Phúc Hậu]
	Đo kích thước các quả đậu Hà Lan ta thu được kết quả:
	\begin{center}
		\begin{tabular}{|c|c|c|c|c|c|c|c|c|c|}
			\hline
			Kích thước & $111$ & $112$ & $113$ & $114$ & $115$ & $116$ & $117$ & $118$ & $119$ \\
			\hline
			Số quả & $3$ & $8$ & $30$ & $68$ & $81$ & $36$ & $18$ & $5$ & $1$ \\
			\hline
		\end{tabular}
	\end{center}
	Tính phương sai của mẫu số liệu.
	\choice
	{\True $1{,}82$}
	{$1{,}71$}
	{$2{,}12$}
	{$1{,}07$}
\loigiai{
	Cỡ mẫu là $N = 3+8+30+68+81+36+18+5+1 = 250$.\\
	Số trung bình cộng của mẫu số liệu là:
	$$\bar{x} = \dfrac{111\cdot 3 + 112\cdot 8 + 113\cdot 30 + \dots + 119\cdot 1}{250} = \dfrac{28677}{250} = 114{,}708$$
	Phương sai của mẫu số liệu là:
	$$s^2 = \dfrac{1}{250}\left[3(111-114{,}708)^2 + 8(112-114{,}708)^2 + \dots + 1(119-114{,}708)^2\right] \approx 1{,}82$$
	}
\end{ex}
\Closesolutionfile{ans}
\indapan{6}{ans/ans-c1b1-lc}
%---------------------Trắc nghiệm đúng sai
\ind{PHẦN II.} \inden{Câu trắc nghiệm đúng sai. Trong mỗi ý a), b), c), d) ở mỗi câu, học sinh chọn đúng hoặc sai.}
\setcounter{ex}{0}
\Opensolutionfile{ans}[ans/ans-c1b1-ds]
\begin{ex}%[0D6H4-2]%[Biên soạn tài liệu dạy học bộ KNTT-CS]%[Phúc Hậu]
	Cuối học kì 1 vừa qua, bạn An đạt được kết quả sáu môn như sau
	\begin{center}
		\begin{tabular}{|c|c|c|c|c|c|c|}
			\hline
			Môn & Toán & Văn & Anh & Lý & Hóa & Sinh \\
			\hline
			Điểm TB & $7{,}2$ & $8{,}0$ & $5{,}8$ & $7{,}2$ & $9{,}0$ & $4{,}6$ \\
			\hline
		\end{tabular}
	\end{center}
	\choiceTF
	{Điểm trung bình các môn thi học kì của bạn An là $7{,}0$}
	{Khoảng biến thiên của bảng điểm của bạn An bằng $3{,}4$}
	{\True Trung vị của bảng điểm của bạn An bằng $7{,}2$}
	{\True Khoảng tứ phân vị bảng điểm của bạn An bằng $2{,}2$}
\loigiai{
	Sắp xếp mẫu số liệu theo thứ tự không giảm: $4{,}6;\ 5{,}8;\ 7{,}2;\ 7{,}2;\ 8{,}0;\ 9{,}0$.
	\begin{itemize}
		\item Số trung bình cộng: $\bar{x} = \dfrac{4{,}6+5{,}8+7{,}2+7{,}2+8{,}0+9{,}0}{6} = \dfrac{41{,}8}{6} \approx 6{,}97$. Vậy ý a) \textbf{Sai}.
		\item Khoảng biến thiên: $R = 9{,}0 - 4{,}6 = 4{,}4$. Vậy ý b) \textbf{Sai}.
		\item Cỡ mẫu $n=6$ (chẵn), trung vị là trung bình cộng hai số ở giữa: $Q_2 = \dfrac{7{,}2+7{,}2}{2} = 7{,}2$. Vậy ý c) \textbf{Đúng}.
		\item Tứ phân vị thứ nhất (trung vị của $4{,}6; 5{,}8; 7{,}2$) là $Q_1 = 5{,}8$.
		\item Tứ phân vị thứ ba (trung vị của $7{,}2; 8{,}0; 9{,}0$) là $Q_3 = 8{,}0$.
		\item Khoảng tứ phân vị $\Delta_Q = Q_3 - Q_1 = 8{,}0 - 5{,}8 = 2{,}2$. Vậy ý d) \textbf{Đúng}.
	\end{itemize}
	}
\end{ex}
%2
\begin{ex}%[0D6H4-4]%[Biên soạn tài liệu dạy học bộ KNTT-CS]%[Phúc Hậu]
	Hai xạ thủ A và B mỗi xạ thủ bắn $10$ phát đạn. Kết quả được thể hiện trong bảng sau
	\begin{center}
		\begin{tabular}{|c|c|c|c|c|c|c|c|c|c|c|}
			\hline
			Xạ thủ A & $7$ & $9$ & $6$ & $9$ & $8$ & $6$ & $8$ & $7$ & $10$ & $8$ \\
			\hline
			Xạ thủ B & $8$ & $7$ & $8$ & $9$ & $6$ & $7$ & $7$ & $9$ & $9$ & $8$ \\
			\hline
		\end{tabular}
	\end{center}
	\choiceTF
	{\True Điểm thấp nhất của xạ thủ A là $6$}
	{Điểm trung bình của xạ thủ A cao hơn điểm trung bình của xạ thủ B}
	{\True Độ lệch chuẩn bảng điểm của xạ thủ A lớn hơn độ lệch chuẩn bảng điểm của xạ thủ B}
	{Xạ thủ A bắn đều hơn xạ thủ B}
\loigiai{
	\begin{itemize}
		\item Quan sát bảng số liệu, điểm thấp nhất của xạ thủ A là $6$. Vậy ý a) \textbf{Đúng}.
		\item Điểm trung bình của xạ thủ A: $\bar{x}_A = \dfrac{7+9+6+9+8+6+8+7+10+8}{10} = 7{,}8$.
		\item Điểm trung bình của xạ thủ B: $\bar{x}_B = \dfrac{8+7+8+9+6+7+7+9+9+8}{10} = 7{,}8$.
		\item Vì $\bar{x}_A = \bar{x}_B$ nên ý b) \textbf{Sai}.
		\item Phương sai của xạ thủ A: $s^2_A = \dfrac{1}{10}\left[2(6-7{,}8)^2 + 2(7-7{,}8)^2 + 3(8-7{,}8)^2 + 2(9-7{,}8)^2 + 1(10-7{,}8)^2\right] = 1{,}56$.
		\item Phương sai của xạ thủ B: $s^2_B = \dfrac{1}{10}\left[1(6-7{,}8)^2 + 3(7-7{,}8)^2 + 3(8-7{,}8)^2 + 3(9-7{,}8)^2\right] = 0{,}96$.
		\item Vì $s^2_A > s^2_B \Rightarrow s_A > s_B$. Vậy ý c) \textbf{Đúng}.
		\item Vì phương sai (độ lệch chuẩn) càng nhỏ thì mức độ phân tán càng thấp (bắn càng đều). Do $s_B < s_A$ nên xạ thủ B bắn đều hơn xạ thủ A. Vậy ý d) \textbf{Sai}.
	\end{itemize}
	}
\end{ex}
%3
\begin{ex}%[0D6H4-4]%[Biên soạn tài liệu dạy học bộ KNTT-CS]%[Phúc Hậu]
	Mẫu số liệu sau đây cho biết chiều cao của $10$ học sinh (đơn vị cm)
	$$165 \quad 155 \quad 160 \quad 145 \quad 157 \quad 162 \quad 148 \quad 170 \quad 172 \quad 152$$
	\choiceTF
	{Chiều cao trung bình của $10$ học sinh là $157{,}6$}
	{\True Khoảng biến thiên của mẫu số liệu là $27$}
	{\True Khoảng tứ phân vị của mẫu số liệu là $13$}
	{Phương sai của mẫu số liệu trên nhỏ hơn $64$}
\loigiai{
	Sắp xếp mẫu số liệu: $145;\ 148;\ 152;\ 155;\ 157;\ 160;\ 162;\ 165;\ 170;\ 172$.
	\begin{itemize}
		\item Chiều cao trung bình: $\bar{x} = \dfrac{145+148+152+155+157+160+162+165+170+172}{10} = 158{,}6$. Vậy ý a) \textbf{Sai}.
		\item Khoảng biến thiên: $R = 172 - 145 = 27$. Vậy ý b) \textbf{Đúng}.
		\item Tứ phân vị thứ nhất $Q_1$ (trung vị của $145; 148; 152; 155; 157$) là $152$.
		\item Tứ phân vị thứ ba $Q_3$ (trung vị của $160; 162; 165; 170; 172$) là $165$.
		\item Khoảng tứ phân vị $\Delta_Q = 165 - 152 = 13$. Vậy ý c) \textbf{Đúng}.
		\item Phương sai mẫu số liệu: $s^2 \approx 72{,}04 > 64$. Vậy ý d) \textbf{Sai}.
	\end{itemize}
	}
\end{ex}
%4
\begin{ex}%[0D6H4-4]%[Biên soạn tài liệu dạy học bộ KNTT-CS]%[Phúc Hậu]
	Dưới đây là điểm kiểm tra giữa kì I của hai bạn An và Bình:
	\begin{center}
		\begin{tabular}{|c|c|c|c|c|c|c|c|c|c|c|c|}
			\hline
			& Toán & Lý & Hóa & Sinh & Anh & Văn & Sử & Địa & CD & Tin & CN \\
			\hline
			An & $8$ & $7{,}5$ & $7{,}8$ & $8{,}3$ & $9$ & $7$ & $8$ & $8{,}2$ & $9$ & $8$ & $8{,}3$ \\
			\hline
			Bình & $8{,}5$ & $9{,}5$ & $9{,}5$ & $8{,}5$ & $9$ & $5$ & $5{,}5$ & $6$ & $10$ & $9$ & $8{,}5$ \\
			\hline
		\end{tabular}
	\end{center}
	\choiceTF
	{\True Khoảng biến thiên của mẫu số liệu điểm kiểm tra của An là $2$}
	{Điểm kiểm tra trung bình của học sinh An là $8{,}3$}
	{\True Điểm kiểm tra trung bình của học sinh Bình là $8{,}1$}
	{\True Bạn An học đều hơn bạn Bình}
\loigiai{
	Số môn học là $n=11$.
	\begin{itemize}
		\item \textbf{Xét kết quả của bạn An:}
		Sắp xếp điểm số theo thứ tự không giảm:
		$$7; \ 7{,}5; \ 7{,}8; \ 8; \ 8; \ 8; \ 8{,}2; \ 8{,}3; \ 8{,}3; \ 9; \ 9$$
		Giá trị lớn nhất là $9$, giá trị nhỏ nhất là $7$.
		Khoảng biến thiên: $R_{An} = 9 - 7 = 2$.
		$\Rightarrow$ Ý a) \textbf{Đúng}.
		Điểm trung bình của An:
		$$\bar{x}_{An} = \dfrac{8+7{,}5+7{,}8+8{,}3+9+7+8+8{,}2+9+8+8{,}3}{11} = \dfrac{89{,}1}{11} = 8{,}1$$
		$\Rightarrow$ Ý b) \textbf{Sai} (vì $8{,}1 \neq 8{,}3$).
		\item \textbf{Xét kết quả của bạn Bình:}
		Tổng điểm của Bình:
		$$\Sigma x_{Bình} = 8{,}5+9{,}5+9{,}5+8{,}5+9+5+5{,}5+6+10+9+8{,}5 = 89$$
		Điểm trung bình của Bình:
		$$\bar{x}_{Bình} = \dfrac{89}{11} \approx 8{,}09$$
		Làm tròn đến hàng phần mười ta được $8{,}1$.
		$\Rightarrow$ Ý c) \textbf{Đúng}.
		\item \textbf{So sánh độ ổn định (học đều):}
		Khoảng biến thiên của Bình là $R_{Bình} = 10 - 5 = 5$.
		Ta thấy $R_{An} = 2 < R_{Bình} = 5$. Đồng thời, các điểm số của An tập trung sát giá trị trung bình ($7$ đến $9$), trong khi điểm của Bình phân tán rộng ($5$ đến $10$).
		Điều này cho thấy phương sai (độ lệch chuẩn) của An nhỏ hơn của Bình, tức là An có kết quả ổn định hơn (học đều hơn).
		$\Rightarrow$ Ý d) \textbf{Đúng}.
	\end{itemize}
	}
\end{ex}
\Closesolutionfile{ans}
\indapan{3}{ans/ans-c1b1-ds}
%---------------Trắc nghiệm trả lời ngắn
\ind{PHẦN III.} \inden{Câu trả lời ngắn.}
\setcounter{ex}{0}
\Opensolutionfile{ans}[ans/ans-c1b1-kq]
%-----Câu 1
\begin{ex}%[0D6N4-2]%[Biên soạn tài liệu dạy học bộ KNTT-CS]%[Phúc Hậu]
	Kết quả điểm thi môn Toán của Tổ $1$, Lớp $10A1$ cho bằng mẫu số liệu:
	$8$; $3$; $5$; $9$; $10$; $9$; $7$; $8$; $2$; $9$; $4$. Khoảng biến thiên của mẫu số liệu bằng?
	\par
	\shortans{$8$}
\loigiai{
	Sắp xếp mẫu số liệu theo thứ tự không giảm: $2$; $3$; $4$; $5$; $7$; $8$; $8$; $9$; $9$; $9$; $10$.\\
	Giá trị lớn nhất của mẫu số liệu là $10$, giá trị nhỏ nhất là $2$.\\
	Khoảng biến thiên của mẫu số liệu là $R=10-2=8$.
	}
\end{ex}
%2
\begin{ex}%[0D6H4-4]%[Biên soạn tài liệu dạy học bộ KNTT-CS]%[Phúc Hậu]
	Thời gian chạy cự li $100$m của các học sinh lớp $10A1$ được cho bằng bảng sau
	\begin{center}
		\begin{tabular}{|c|c|c|c|c|c|}
			\hline
			Thời gian & $12$ & $13$ & $14$ & $15$ & $16$ \\
			\hline
			Số bạn    & $5$  & $7$  & $10$ & $8$  & $6$  \\
			\hline
		\end{tabular}
	\end{center}
	Tính phương sai của mẫu số liệu trên (làm tròn đến hàng phần chục).
	\par
	\shortans{$1{,}6$}
\loigiai{
	Số trung bình cộng của mẫu số liệu là
	\[
	\overline{x} = \dfrac{12\cdot 5 + 13\cdot 7 + 14\cdot 10 + 15\cdot 8 + 16\cdot 6}{5+7+10+8+6} = \dfrac{507}{36} \approx 14{,}08.
	\]
	Phương sai của mẫu số liệu là
	\[
	s^2 = \dfrac{1}{36}\left[5\cdot(12-14{,}08)^2 + 7\cdot(13-14{,}08)^2 + 10\cdot(14-14{,}08)^2 + 8\cdot(15-14{,}08)^2 + 6\cdot(16-14{,}08)^2\right] \approx 1{,}6.
	\]
	Vậy phương sai của mẫu số liệu xấp xỉ $1{,}6$.
	}
\end{ex}
%3
\begin{ex}%[0D6H4-2]%[Biên soạn tài liệu dạy học bộ KNTT-CS]%[Phúc Hậu]
	Kết quả điểm thi môn Toán của Tổ $1$, Lớp $10A1$ cho bằng mẫu số liệu:
	$8$; $3$; $5$; $9$; $10$; $9$; $7$; $8$; $2$; $9$; $4$. Khoảng biến thiên của mẫu số liệu bằng?
	\par
	\shortans{$8$}
\loigiai{
	Sắp xếp mẫu số liệu theo thứ tự không giảm: $2$; $3$; $4$; $5$; $7$; $8$; $8$; $9$; $9$; $9$; $10$.\\
	Giá trị lớn nhất của mẫu số liệu là $10$, giá trị nhỏ nhất là $2$.\\
	Khoảng biến thiên của mẫu số liệu là $R=10-2=8$.
	}
\end{ex}
%4
\begin{ex}%[0D6H4-4]%[Biên soạn tài liệu dạy học bộ KNTT-CS]%[Phúc Hậu]
	Thời gian chạy cự li $100$m của các học sinh lớp $10A1$ được cho bằng bảng sau
	\begin{center}
		\begin{tabular}{|c|c|c|c|c|c|}
			\hline
			Thời gian & $12$ & $13$ & $14$ & $15$ & $16$ \\
			\hline
			Số bạn    & $5$  & $7$  & $10$ & $8$  & $6$  \\
			\hline
		\end{tabular}
	\end{center}
	Tính phương sai của mẫu số liệu trên (làm tròn đến hàng phần chục).
	\par
	\shortans{$1{,}6$}
\loigiai{
	Số trung bình cộng của mẫu số liệu là
	\[
	\overline{x} = \dfrac{12\cdot 5 + 13\cdot 7 + 14\cdot 10 + 15\cdot 8 + 16\cdot 6}{5+7+10+8+6} = \dfrac{507}{36} \approx 14{,}08.
	\]
	Phương sai của mẫu số liệu là
	\[
	s^2 = \dfrac{1}{36}\left[5\cdot(12-14{,}08)^2 + 7\cdot(13-14{,}08)^2 + 10\cdot(14-14{,}08)^2 + 8\cdot(15-14{,}08)^2 + 6\cdot(16-14{,}08)^2\right] \approx 1{,}6.
	\]
	Vậy phương sai của mẫu số liệu xấp xỉ $1{,}6$.
	}
\end{ex}
%5
\begin{ex}%[0D6H4-4]%[Biên soạn tài liệu dạy học bộ KNTT-CS]%[Phúc Hậu]
	Cho mẫu số liệu $\{a; 4; b; 8\}$ với $a < b$. Biết giá trị trung bình của mẫu số liệu $\overline{x} = 5$ và phương sai của mẫu số liệu $s^2 = 5$. Tính $a+b$.
	\par
	\shortans{$8$}
\loigiai{
	Số trung bình của mẫu số liệu là:
	\[ \overline{x} = \dfrac{a + 4 + b + 8}{4} = 5 \Leftrightarrow a + b + 12 = 20 \Leftrightarrow a + b = 8. \]
	Vậy $a+b=8$.
	\textit{Kiểm tra lại (không bắt buộc):}
	Từ phương sai $s^2 = 5$, ta có:
	\[ \dfrac{(a-5)^2 + (4-5)^2 + (b-5)^2 + (8-5)^2}{4} = 5 \Leftrightarrow (a-5)^2 + (b-5)^2 = 10. \]
	Kết hợp với $b=8-a$, giải ra ta được $a=2, b=6$ (thỏa mãn $a<b$).
	}
\end{ex}
%6
\begin{ex}%[0D6H4-3]%[Biên soạn tài liệu dạy học bộ KNTT-CS]%[Phúc Hậu]
	Mẫu số liệu sau đây cho biết cân nặng của $10$ trẻ sơ sinh (đơn vị kg)
	\begin{center}
		\begin{tabular}{|c|c|c|c|c|}
			\hline
			$2{,}977$ & $3{,}155$ & $3{,}920$ & $3{,}412$ & $4{,}236$ \\
			\hline
			$2{,}593$ & $3{,}270$ & $3{,}813$ & $4{,}042$ & $3{,}387$ \\
			\hline
		\end{tabular}
	\end{center}
	Có bao nhiêu giá trị bất thường trong mẫu số liệu đã cho?
	\shortans{$0$}
\loigiai{
	Sắp xếp mẫu số liệu theo thứ tự không giảm:
	\[ 2{,}593;\ 2{,}977;\ 3{,}155;\ 3{,}270;\ 3{,}387;\ 3{,}412;\ 3{,}813;\ 3{,}920;\ 4{,}042;\ 4{,}236 \]
	Cỡ mẫu $n=10$.
	Tứ phân vị thứ nhất $Q_1$ là trung vị của nửa số liệu đầu ($5$ giá trị đầu):
	\[ 2{,}593;\ 2{,}977;\ \mathbf{3{,}155};\ 3{,}270;\ 3{,}387 \Rightarrow Q_1 = 3{,}155. \]
	Tứ phân vị thứ ba $Q_3$ là trung vị của nửa số liệu sau ($5$ giá trị cuối):
	\[ 3{,}412;\ 3{,}813;\ \mathbf{3{,}920};\ 4{,}042;\ 4{,}236 \Rightarrow Q_3 = 3{,}920. \]
	Khoảng tứ phân vị: $\Delta_Q = Q_3 - Q_1 = 3{,}920 - 3{,}155 = 0{,}765$.
	Ta có các giới hạn xác định giá trị bất thường:
	\begin{itemize}
		\item Giới hạn dưới: $Q_1 - 1{,}5\Delta_Q = 3{,}155 - 1{,}5 \cdot 0{,}765 = 2{,}0075$.
		\item Giới hạn trên: $Q_3 + 1{,}5\Delta_Q = 3{,}920 + 1{,}5 \cdot 0{,}765 = 5{,}0675$.
	\end{itemize}
	Quan sát mẫu số liệu, ta thấy không có giá trị nào nhỏ hơn $2{,}0075$ và không có giá trị nào lớn hơn $5{,}0675$.
	Vậy không có giá trị bất thường nào trong mẫu số liệu.
	}
\end{ex}
\Closesolutionfile{ans}
\indapan{4}{ans/ans-c1b1-kq}
%------------------Tự luận
\ind{PHẦN IV.} \inden{Tự luận.}
\setcounter{ex}{0}
%1
\begin{ex}%[0D6H4-4]%[Biên soạn tài liệu dạy học bộ KNTT-CS]%[Phúc Hậu]
	Một công ty nông nghiệp làm thí nghiệm trên $40$ thửa ruộng có cùng diện tích và thu được sản lượng lúa (đơn vị: tạ) được trình bày trong bảng sau:
	\begin{center}
		\begin{tabular}{|l|c|c|c|c|c|}
			\hline
			Sản lượng & $20$ & $21$ & $22$ & $23$ & $24$ \\
			\hline
			Số thửa ruộng & $5$ & $8$ & $11$ & $10$ & $6$ \\
			\hline
		\end{tabular}
	\end{center}
	Tính phương sai của mẫu số liệu trên.
\loigiai{
	Số trung bình của mẫu số liệu là
	\[ \overline{x} = \dfrac{20\cdot 5 + 21\cdot 8 + 22\cdot 11 + 23\cdot 10 + 24\cdot 6}{40} = \dfrac{884}{40} = 22{,}1. \]
	Phương sai của mẫu số liệu là:
	\begin{align*}
		s^2 &= \dfrac{1}{40}\left[ 5(20-22{,}1)^2 + 8(21-22{,}1)^2 + 11(22-22{,}1)^2 + 10(23-22{,}1)^2 + 6(24-22{,}1)^2 \right] \\
		&= \dfrac{1}{40}(22{,}05 + 9{,}68 + 0{,}11 + 8{,}1 + 21{,}66) \\
		&= \dfrac{61{,}6}{40} = 1{,}54.
	\end{align*}
	}
\end{ex}
%2
\begin{ex}%[0D6H4-2]%[Biên soạn tài liệu dạy học bộ KNTT-CS]%[Phúc Hậu]
	Cho mẫu số liệu $13,\ 26,\ 21,\ 17,\ 14,\ 13,\ 25,\ 17,\ 17,\ 10,\ 14,\ 9$. Tính khoảng tứ phân vị của mẫu số liệu.
\loigiai{
	Sắp xếp mẫu số liệu theo thứ tự không giảm ($n=12$):
	\[ 9;\ 10;\ 13;\ 13;\ 14;\ 14;\ 17;\ 17;\ 17;\ 21;\ 25;\ 26. \]
	Tứ phân vị thứ nhất $Q_1$ là trung vị của nửa số liệu đầu ($9;\ 10;\ 13;\ 13;\ 14;\ 14$):
	\[ Q_1 = \dfrac{13+13}{2} = 13. \]
	Tứ phân vị thứ ba $Q_3$ là trung vị của nửa số liệu sau ($17;\ 17;\ 17;\ 21;\ 25;\ 26$):
	\[ Q_3 = \dfrac{17+21}{2} = 19. \]
	Khoảng tứ phân vị của mẫu số liệu là: $\Delta_Q = Q_3 - Q_1 = 19 - 13 = 6$.
	}
\end{ex}
%3
\begin{ex}%[0D6H4-4]%[Biên soạn tài liệu dạy học bộ KNTT-CS]%[Phúc Hậu]
	Cho mẫu số liệu $23,\ 20,\ 29,\ 27,\ 24,\ 29,\ 21,\ 23,\ 26,\ 19,\ 23,\ 27,\ 27$. Tính độ lệch chuẩn của mẫu số liệu.
\loigiai{
	Cỡ mẫu $n=13$.
	Số trung bình cộng của mẫu số liệu là:
	\[ \overline{x} = \dfrac{23+20+29+27+24+29+21+23+26+19+23+27+27}{13} = \dfrac{318}{13} \approx 24{,}46. \]
	Phương sai của mẫu số liệu là:
	\[ s^2 = \dfrac{1}{13}(23^2 + 20^2 + \dots + 27^2) - \left(\dfrac{318}{13}\right)^2 = \dfrac{7910}{13} - \dfrac{101124}{169} \approx 10{,}09. \]
	Độ lệch chuẩn của mẫu số liệu là
	\[ s = \sqrt{s^2} \approx \sqrt{10{,}09} \approx 3{,}2. \]
	}
\end{ex}
%4
\begin{ex}%[0D6H4-4]%[Biên soạn tài liệu dạy học bộ KNTT-CS]%[Phúc Hậu]
	Cho các số liệu thống kê ghi trong bảng sau:\\
	Chiều cao của các học sinh trong một lớp $10$ (đơn vị: cm)
	\begin{center}
		\begin{tabular}{|l|l|l|l|l|l|l|}
			\hline
			$160$ & $161$ & $180$ & $175$ & $177$ & $178$ & $175$ \\
			\hline
			$180$ & $178$ & $172$ & $161$ & $170$ & $171$ & $170$ \\
			\hline
			$168$ & $168$ & $168$ & $160$ & $165$ & $170$ & $165$ \\
			\hline
			$161$ & $170$ & $164$ & $177$ & $157$ & $160$ & $161$ \\
			\hline
			$165$ & $177$ & $165$ & $178$ & $164$ & $175$ & $157$ \\
			\hline
			$170$ & $172$ & $159$ & $155$ & $159$ & $165$ & $155$ \\
			\hline
		\end{tabular}
	\end{center}
	Tính phương sai của mẫu số liệu trên (làm tròn kết quả đến hàng phần mười).
\loigiai{
	Tổng số học sinh là $N = 42$.
	Ta lập bảng phân bố tần số:
	\begin{center}
		\begin{tabular}{|c|c|c|c|c|c|c|c|c|c|c|c|c|c|c|c|}
			\hline
			$x_i$ & $155$ & $157$ & $159$ & $160$ & $161$ & $164$ & $165$ & $168$ & $170$ & $171$ & $172$ & $175$ & $177$ & $178$ & $180$ \\
			\hline
			$n_i$ & $2$ & $2$ & $2$ & $3$ & $4$ & $2$ & $5$ & $3$ & $5$ & $1$ & $2$ & $3$ & $3$ & $3$ & $2$ \\
			\hline
		\end{tabular}
	\end{center}
	Số trung bình cộng của mẫu số liệu là:
	\[ \overline{x} = \dfrac{1}{42} \sum n_i x_i = \dfrac{7038}{42} = \dfrac{1173}{7} \approx 167{,}57. \]
	Phương sai của mẫu số liệu được tính theo công thức $s^2 = \overline{x^2} - (\overline{x})^2$:
	\[ s^2 = \dfrac{1}{42} \sum n_i x_i^2 - (\overline{x})^2 = \dfrac{1\,181\,606}{42} - \left(\dfrac{1173}{7}\right)^2 \approx 53{,}29. \]
	Làm tròn kết quả đến hàng phần mười, phương sai là $53{,}3$.
	}
\end{ex}
%5
\begin{ex}%[0D6H4-4]%[Biên soạn tài liệu dạy học bộ KNTT-CS]%[Phúc Hậu]
	Điểm thi học sinh trong đợt kiểm tra đầu vào môn Tiếng Anh (thang điểm $100$) của hai lớp $3A$ và $3B$ được xác định bởi bảng sau:
	\begin{center}
		\begin{tabular}{|l|c|c|c|c|c|c|c|c|c|c|c|c|}
			\hline
			Lớp 3A & $45$ & $50$ & $55$ & $60$ & $65$ & $70$ & $75$ & $80$ & $85$ & $90$ & $95$ & $100$ \\
			\hline
			Tần số & $2$ & $3$ & $3$ & $4$ & $5$ & $8$ & $6$ & $3$ & $0$ & $2$ & $3$ & $1$ \\
			\hline
		\end{tabular}
	\end{center}
	\begin{center}
		\begin{tabular}{|l|c|c|c|c|c|c|c|c|c|c|c|c|}
			\hline
			Lớp 3B & $45$ & $50$ & $55$ & $60$ & $65$ & $70$ & $75$ & $80$ & $85$ & $90$ & $95$ & $100$ \\
			\hline
			Tần số & $0$ & $7$ & $8$ & $4$ & $1$ & $0$ & $1$ & $3$ & $3$ & $4$ & $4$ & $3$ \\
			\hline
		\end{tabular}
	\end{center}
	Tính tỉ số phương sai của học sinh $3B$ với học sinh $3A$ (Kết quả làm tròn đến hàng phần trăm).
\loigiai{
	\textbf{Đối với lớp 3A:}
	Tổng số học sinh là $N_A = 2+3+3+4+5+8+6+3+0+2+3+1 = 40$.
	Số trung bình cộng điểm thi của lớp $3A$ là:
	\[ \overline{x}_A = \dfrac{45\cdot 2 + 50\cdot 3 + \dots + 100\cdot 1}{40} = \dfrac{2785}{40} = 69{,}625. \]
	Phương sai điểm thi của lớp $3A$ là:
	\[ s^2_A = \dfrac{1}{40}(2\cdot 45^2 + 3\cdot 50^2 + \dots + 1\cdot 100^2) - (69{,}625)^2 = 191{,}734375. \]
	\textbf{Đối với lớp 3B:}
	Tổng số học sinh là $N_B = 0+7+8+4+1+0+1+3+3+4+4+3 = 38$.
	Số trung bình cộng điểm thi của lớp $3B$ là:
	\[ \overline{x}_B = \dfrac{45\cdot 0 + 50\cdot 7 + \dots + 100\cdot 3}{38} = \dfrac{2705}{38} \approx 71{,}18. \]
	Phương sai điểm thi của lớp $3B$ là:
	\[ s^2_B = \dfrac{1}{38}(0\cdot 45^2 + 7\cdot 50^2 + \dots + 3\cdot 100^2) - \left(\dfrac{2705}{38}\right)^2 \approx 336{,}10. \]
	\textbf{Tỉ số phương sai:}
	Tỉ số phương sai của học sinh $3B$ so với học sinh $3A$ là:
	\[ \dfrac{s^2_B}{s^2_A} \approx \dfrac{336{,}10}{191{,}73} \approx 1{,}75. \]
	}
\end{ex}
%6
\begin{ex}%[0D6H4-3]%[Biên soạn tài liệu dạy học bộ KNTT-CS]%[Phúc Hậu]
	Một cảnh sát giao thông ghi tốc độ (đơn vị km/h) của $25$ chiếc xe qua trạm như sau:
	\begin{center}
		\begin{tabular}{ccccccccccccc}
			$20$ & $41$ & $41$ & $80$ & $40$ & $52$ & $52$ & $52$ & $60$ & $55$ & $60$ & $60$ & $62$ \\
			$60$ & $65$ & $60$ & $65$ & $135$ & $70$ & $70$ & $65$ & $75$ & $75$ & $70$ & $55$ &
		\end{tabular}
	\end{center}
	Tìm giá trị bất thường trong bảng số liệu trên.
\loigiai{
	Sắp xếp mẫu số liệu theo thứ tự không giảm:
	\[ 20;\ 40;\ 41;\ 41;\ 52;\ 52;\ 52;\ 55;\ 55;\ 60;\ 60;\ 60;\ 60;\ 60;\ 62;\ 65;\ 65;\ 65;\ 70;\ 70;\ 70;\ 75;\ 75;\ 80;\ 135. \]
	Cỡ mẫu $n=25$.
	Trung vị $Q_2$ là giá trị ở vị trí thứ $13$, suy ra $Q_2 = 60$.
	Tứ phân vị thứ nhất $Q_1$ là trung vị của nửa số liệu bên trái $Q_2$ (gồm $12$ giá trị đầu):
	\[ 20;\ 40;\ 41;\ 41;\ 52;\ 52;\ 52;\ 55;\ 55;\ 60;\ 60;\ 60. \]
	Suy ra $Q_1 = \dfrac{52+52}{2} = 52$.
	Tứ phân vị thứ ba $Q_3$ là trung vị của nửa số liệu bên phải $Q_2$ (gồm $12$ giá trị cuối):
	\[ 60;\ 62;\ 65;\ 65;\ 65;\ 70;\ 70;\ 70;\ 75;\ 75;\ 80;\ 135. \]
	Suy ra $Q_3 = \dfrac{70+70}{2} = 70$.
	Khoảng tứ phân vị của mẫu số liệu là:
	\[ \Delta_Q = Q_3 - Q_1 = 70 - 52 = 18. \]
	Ta xác định các giới hạn để tìm giá trị bất thường:
	\begin{itemize}
		\item Giới hạn dưới: $Q_1 - 1{,}5\Delta_Q = 52 - 1{,}5 \cdot 18 = 25$.
		\item Giới hạn trên: $Q_3 + 1{,}5\Delta_Q = 70 + 1{,}5 \cdot 18 = 97$.
	\end{itemize}
	Trong mẫu số liệu, ta thấy:
	\begin{itemize}
		\item Giá trị $20 < 25$.
		\item Giá trị $135 > 97$.
	\end{itemize}
	Vậy các giá trị bất thường trong mẫu số liệu là $20$ và $135$.
	}
\end{ex}
\Closesolutionfile{ans}